\documentclass{article}
\usepackage{amsmath} %for align, etc.
\usepackage{amssymb} %for \mathbb, \mathcal, etc.
\usepackage{graphicx} %for resizing large displayed matrices
\usepackage{comment} %to comment out some stuff we might use later instead
\newcommand{\lbar}[1]{\underline{#1}} %easier way to write lower bar
\DeclareMathOperator\supp{supp} %support function

\begin{document}
\textit{To solve this problem, we first make the assumption that all young workers initially join the labour-force unemployed. We also assume that unemployed workers who search can only accept their offer in the next period.}

\section{A General Analytic Solution}
\smallskip
Since old workers leave the labour-force upon retirement, this amounts to all workers facing the following age transition matrix:\footnote{We have opted to replace all $\pi$s in the age transition probabilities with $\gamma$s so as to keep this notationally distinct from the endogenous job finding probability $\pi(a)$}
\[
\Gamma = \begin{bmatrix}
    \gamma_{jj} & 1-\gamma_{jj} & 0 \\
    0 & \gamma_{mm} & 1-\gamma_{mm} \\
    0 & 0 & \gamma_{ss}
\end{bmatrix}
\]
In effect, in each period of their existence, senior workers exit the labour-force with a probability of $1-\gamma_{ss}$. As we will see later, this places heterogeneous incentives on the search effort of unemployed workers throughout their life cycle.

Let $V_{l,g}$ denote the value function of a worker who is of employment status $l\in\{e,u\}$ (employed or unemployed), and age $g\in\{j,m,s\}$ (young, middle-aged, or senior).\footnote{Note that we have opted to use $g$ to index age, to prevent any possible future confusion with the senior age $s$, accordingly, we have also used $\delta$ for the exogenous separation/job destruction rate.} At each period, a worker faces a stochastic infinite horizon problem, where each subsequent period they experience a probability of aging up (or leaving the labour-force if senior), and changing employment status. Hence, we write the representative Bellman equations as follows, for all ages $g\in\{j,m,s\}$:
\begin{align}
V_{e,g}(w)
&=
w+\beta\sum_{g'}\Gamma(g'\mid g)
\Bigl[
(1-\delta)V_{e,g'}(w)+\delta V_{u,g'}
\Bigr],
\tag{W}\label{eq:bell_w}
\\
V_{u,g}
&= \max_{a_g \ge 0} \Biggl\{
b - \psi a_g + \beta \sum_{g'} \Gamma(g' \mid g)
\Bigl[
(1 - \pi(a_g)) V_{u,g'} \tag{U}~\label{eq:bell_u}\\
&\qquad\qquad + \pi(a_g)\, \mathbb{E}_w \bigl[\max\{V_{u,g'}, V_{e,g'}(w)\}\bigr]
\Bigr]
\Biggr\} \notag
\end{align}
Here, we let $g$ refer to the age of worker in their present period, and $g'$ to their age in their next period. Similarly, we use $\Gamma(g'|g)$ to refer to the age transition probability to age $g'$ given they are currently at age $g$, where combining with the sum over all $g'$ gives the expected next period age. Altogether,~\eqref{eq:bell_w} and~\eqref{eq:bell_u} represent the Bellman equations for employed and unemployed workers respectively representing the present values of each state, where $V_{e,g}(w)$ takes the current wage as an argument through its role as a state variable.

We pay particular attention to~\eqref{eq:bell_u}, insofar as the wage itself is drawn stochastically from $F(w)$. Hence, a worker only accepts an offer if the present value of a accepting the offer outweighs that of being unemployed, represented in the ${E}_w \bigl[\max\{V_{u,g'}, V_{e,g'}(w)\}\bigr]$ portion of the present value of being unemployed.

Solving backwards from the final age cohort, we note that if a senior worker is in a given labour market state in the next period, they can only be so when senior. With this in mind, we obtain the following expression for the value of being 
employed with a given wage $w$ when old:
\begin{align*}
V_{e,s}(w)&=w+\beta\gamma_{ss}\bigl[(1-\delta)V_{e,s}(w)+\delta V_{u,s}\bigr]\\
&=\frac{w+\beta\gamma_{ss}\delta V_{u,s}}{1-\beta\gamma_{ss}(1-\delta)}
\end{align*}
Furthermore, we note that a senior worker will only accept a job if the net present value of accepting the job is at least as great as that of being unemployed:
\begin{align*}
V_{u,s}&\leq\frac{w+\beta\gamma_{ss}\delta V_{u,s}}{1-\beta\gamma_{ss}(1-\delta)}\\
w&\geq(1-\beta\gamma_{ss})V_{u,s}
\end{align*}
Since $F(w)$ is continuous and has full support on $\mathbb{R}$, there exists a reservation wage $\lbar{w}_s\in\supp(F(w))$ satisfying the above expression with perfect equality, such that any wage offers below this value will not be accepted when senior. In effect, a senior worker values unemployment at:
\begin{equation}
V_{u,s}=\frac{\lbar{w}_s}{1-\beta\gamma_{ss}}
~\label{eq:sen_res_unemp}
\end{equation}
Moreover, to compute the expected gain from accepting an offer at a given wage $w\geq\lbar{w}_s$, we can calculate the difference for senior workers between the value of unemployment and that of employment:
\begin{align*}
\mathbb{E}[V_{e,s}(w)-V_{u,s}|w\geq\lbar{w}_s]=\frac{\mathbb{E}[w-\lbar{w}_s|w\geq\lbar{w}_s]}{1-\beta\gamma_{ss}(1-\delta)}\\
\mathbb{E}[V_{e,s}(w)|w\geq\lbar{w}_s]=\frac{\mathbb{E}[w-\lbar{w}_s|w\geq\lbar{w}_s]}{1-\beta\gamma_{ss}(1-\delta)}+V_{u,s}
\end{align*}
Substituting this identity back into the senior's Bellman equation when unemployed, we get:
\begin{align*}
V_{u,s}&=\max_{a_s\geq 0}\Biggl\{b-\psi a_s+\beta\gamma_{ss}\biggl[V_{u,s}+\frac{\pi(a_s)}{1-\beta\gamma_{ss}(1-\delta)}\int_{\lbar{w}_s}^\infty(w-\lbar{w}_s)dF(w)\biggr]\Biggr\}\\
\lbar{w}_s&=\max_{a_s\geq 0}\Biggl\{b-\psi a_s+\frac{\beta\gamma_{ss}\pi(a_s)}{1-\beta\gamma_{ss}(1-\delta)}\int_{\lbar{w}_s}^\infty(w-\lbar{w}_s)dF(w)\Biggr\}\\
\end{align*}
where in the first line we replace the expectation term with the integral with respect to the distribution, and in the second regroup $V_{u,s}$ to the left-hand-side, and substitute for the senior reservation wage, given the identity~\eqref{eq:sen_res_unemp}.

To obtain an implicit solution, we first solve for optimal senior search effort $a_s$, taking the first order condition of our expression for $\lbar{w}_s$ with respect to this, we get:
\[
a_s^*=\max\Biggl\{0,\frac{1}{\phi}\ln\biggl(\frac{\phi\beta\gamma_{ss}\int_{\lbar{w}_s}^\infty(w-\lbar{w}_s)dF(w)}{\psi(1-\beta\gamma_{ss}(1-\delta))}\biggr)\Biggr\},
\]
where an interior solution follows only if $\phi\beta\gamma_{ss}\int_{\lbar{w}_s}^\infty(w-\lbar{w}_s)dF(w)>\psi(1-\beta\gamma_{ss}(1-\delta))$, else the senior worker opts to remain unemployed throughout, and the distribution of wages for senior workers is degenerate.

Since $F(w)$ is distributed normally, there is no elementary closed form for the cumulative distribution function. Nonetheless, it is clear that plugging $a_s^*$ back in will provide us with a fixed point equation, hence a there exists a solution to our equation. Moreover, note by Tonelli's theorem that:
\[
\int_{\lbar{w}_s}^{\infty} (w-\lbar{w}_s)dF(w)
= \int_{\lbar{w}_s}^{\infty}\int_x^{\infty} dF(w)dx
= \int_{\lbar{w}_s}^{\infty} \bigl(1-F(w)\bigr)dw.
\]
We can hence write our fixed-point equation for $\lbar{w}_s$ as:
\begin{equation}
\lbar{w}_s=b-\psi a_s^*+\frac{\beta\gamma_{ss}(1-e^{-\phi a_s^*})}{1-\beta\gamma_{ss}(1-\delta)}\int_{\lbar{w}_s}^{\infty} \bigl(1-F(w)\bigr)dw
~\label{eq:ws_fixed_point}
\end{equation}
Intuitively, the right-hand side captures $b$ denotes the flow payoff from unemployment net of the optimal flow search effort cost $\psi a_s^*$, plus the expected wage gain from future job offers: the probability of receiving an acceptable offer, $1-e^{-\phi a_s^*}$, multiplied by the discounted expected surplus from such an offer, while the factor $\frac{\beta\gamma_{ss}}{1-\beta\gamma_{ss}(1-\delta)}$ scales this surplus by the effective duration of employment, accounting for both discounting and job destruction.

\noindent It follows that, at an interior solution:
\[
\frac{d}{d\lbar{w}_s}\biggl(b-\psi a_s^*+\frac{\beta\gamma_{ss}(1-e^{-\phi a_s^*})}{1-\beta\gamma_{ss}(1-\delta)}\int_{\lbar{w}_s}^{\infty} \bigl(1-F(w)\bigr)dw\biggr)<0
\]
To see this explicitly, let us denote the right hand side of~\eqref{eq:ws_fixed_point} by $\lbar{w}_s$, which is itself a function both directly of $\lbar{w}_s$, and implicitly by $a^*_g(\lbar{w}_s)$. Thus, by the envelope theorem:
\begin{align*}
\frac{d\lbar{w}_s(a^*_g(\lbar{w}_s), \lbar{w}_s)}{d\lbar{w}_s}&=\frac{\partial \lbar{w}_s(a^*_g(\lbar{w}_s), \lbar{w}_s)}{\partial a^*_g(\lbar{w}_s)}\frac{\partial a^*_g(\lbar{w}_s)}{\partial \lbar{w}_s}+\frac{\partial \lbar{w}_s(a^*_g(\lbar{w}_s), \lbar{w}_s)}{\partial \lbar{w}_s}\\
&=\frac{\partial \lbar{w}_s(a^*_g(\lbar{w}_s), \lbar{w}_s)}{\partial \lbar{w}_s}\\
&=-\frac{\beta\gamma_{ss}(1-e^{-\phi a_s^*})}{1-\beta\gamma_{ss}(1-\delta)}\bigl(1-F(\lbar{w}_s)\bigr)<0
\end{align*}
Therefore, since $\lbar{w}_s$ is strictly decreasing in $\lbar{w}_s$, and the identity function is strictly increasing, it follows that the senior reservation wage is unique at the interior. At the corner solution, the reservation wage is trivially unique; with $a^*_s=0$, the identity simplifies to:
\[
\lbar{w}_s=b
\]
More importantly, this also affords us the ability to derive implicit solutions for the remaining two age cohorts. With our senior reservation wage now well defined, we proceed onto the subsequent age cohorts. In particular, for notational simplicity, let us define from our solved objects, the present-discounted transition value of being senior for an employed middle-aged worker:
\[
\tilde{V}_s(w)\doteq \beta(1-\gamma_{mm})\bigl((1-\delta)V_{e,s}(w)+\delta V_{u,s}\bigr)
\]
It follows that the middle-aged present value of being employed at wage $w$, as a function of the value of being unemployed is:
\[
V_{e,m}(w)=\frac{w+\beta\gamma_{mm}\delta V_{u,m}+\tilde{V}_s(w)}{1-\beta\gamma_{mm}(1-\delta)}
\]
Proceeding as before, we note the existence of a reservation wage satisfying:
\begin{align*}
\lbar{w}_m=(1-\beta\gamma_{mm})V_{u,m}-\tilde{V}_s(\lbar{w}_m)
\end{align*}
Hence, using this reservation wage condition, we have:
\[
V_{e,m}(w)-V_{u,m}=
\frac{
w-\lbar{w}_m+\tilde V_s(w)-\tilde V_s(\lbar{w}_m)
}
{1-\beta\gamma_{mm}(1-\delta)}
\]
In addition, using the same approach as before, solving for $\lbar{w}_m$, we get:
\begin{align*}
\lbar{w}_m
=
\max_{a_m\geq 0}
\Biggl\{
& b-\psi a_m
+\tilde V_s(\lbar w_s)-\tilde V_s(\lbar w_m) \\
&\quad
+\frac{\beta\gamma_{mm}\pi(a_m)}
{1-\beta\gamma_{mm}(1-\delta)}
\int_{\lbar{w}_m}^{\infty}
\Bigl[
w-\lbar{w}_m+\tilde V_s(w)-\tilde V_s(\lbar{w}_m)
\Bigr]\,dF(w) \\
&\quad
+\frac{\beta(1-\gamma_{mm})\pi(a_m)}
{1-\beta\gamma_{ss}(1-\delta)}
\int_{\lbar{w}_s}^{\infty}
(w-\lbar{w}_s)\,dF(w)
\Biggr\}
\end{align*}
where we use the fact that $V_{e,s}(\lbar{w}_s)=V_{u,s}$. Moreover, this entails an optimal search effort of:
\begin{align*}
a_m^*
=
\max\Biggl\{
0,\,
\frac{1}{\phi}
\ln\Biggl(\frac{\phi}{\psi}\biggl[
&\frac{
\beta\gamma_{mm}}
{1-\beta\gamma_{mm}(1-\delta)}
\int_{\lbar{w}_m}^{\infty}
\Bigl[
w-\lbar{w}_m+\tilde V_s(w)-\tilde V_s(\lbar{w}_m)
\Bigr]\,dF(w)\\
&\quad
+\frac{\beta(1-\gamma_{mm})}
{1-\beta\gamma_{ss}(1-\delta)}
\int_{\lbar{w}_s}^{\infty}
(w-\lbar{w}_s)\,dF(w)
\biggr]\Biggr)
\Biggr\}.
\end{align*}
In particular, the last term captures within it the additional gain when senior, if a suitable wage was drawn when middle-aged:
\begin{align*}
\lbar w_s=b-\psi a_s^*+
\frac{\beta\gamma_{ss}\pi(a_s^*)}
{1-\beta\gamma_{ss}(1-\delta)}
\int_{\lbar w_s}^{\infty}(w-\lbar w_s)\,dF(w)\\
\frac{1}{1-\beta\gamma_{ss}(1-\delta)}
\int_{\lbar w_s}^{\infty}(w-\lbar w_s)\,dF(w)
=
\frac{\lbar w_s-b+\psi a_s^*}
{\beta\gamma_{ss}\pi(a_s^*)}.
\end{align*}
However, while the middle-aged worker inherits a guaranteed portion of the senior worker surplus per the probability weighting of their problem, the relative search effort remains ambiguous. Altogether, we can write the fixed-point equation for $\lbar{w}_m$ as:
\begin{align*}
\lbar{w}_m
=
& \;b-\psi a^*_m
+\tilde V_s(\lbar w_s)-\tilde V_s(\lbar w_m) \\
&\quad
+\frac{\beta\gamma_{mm}\pi(a^*_m)}
{1-\beta\gamma_{mm}(1-\delta)}
\int_{\lbar{w}_m}^{\infty}
\Bigl[
w-\lbar{w}_m+\tilde V_s(w)-\tilde V_s(\lbar{w}_m)
\Bigr]\,dF(w) \\
&\quad
+\frac{\beta(1-\gamma_{mm})\pi(a^*_m)}
{1-\beta\gamma_{ss}(1-\delta)}
\int_{\lbar{w}_s}^{\infty}
(w-\lbar{w}_s)\,dF(w)
\end{align*}
Once more, the resulting size of the reservation wage of middle-aged worker relative to that of a senior worker remains ambiguous. In particular, due to the increased search effort, this reduces the base net flow value. Furthermore, while the gains of a higher wage push up the potential future horizon payoffs, this also means that jobs at a lower wages have higher payoffs due to their increased time horizon.

Finally, we solve the case of the young worker by drawing upon the symmetry of the recursion. In particular:
\begin{align*}
\lbar{w}_j
=
b-\psi a_j^*
&+\tilde V_m(\lbar w_m)-\tilde V_m(\lbar w_j)\\
&+\frac{\beta\gamma_{jj}\pi(a_j^*)}
{1-\beta\gamma_{jj}(1-\delta)}
\int_{\lbar{w}_j}^{\infty}
\Bigl[
w-\lbar{w}_j+\tilde V_m(w)-\tilde V_m(\lbar{w}_j)
\Bigr]\,dF(w)\\
&+\frac{\beta(1-\gamma_{jj})\pi(a_j^*)}
{1-\beta\gamma_{mm}(1-\delta)}
\int_{\lbar{w}_m}^{\infty}
\Bigl[
w-\lbar{w}_m+\tilde V_s(w)-\tilde V_s(\lbar{w}_m)
\Bigr]\,dF(w)
\end{align*}
\begin{align*}
a_j^*
=
\max\Biggl\{
0,\,
\frac{1}{\phi}&
\ln\Biggl(
\frac{\phi}{\psi}
\biggl[
\frac{\beta\gamma_{jj}}
{1-\beta\gamma_{jj}(1-\delta)}
\int_{\lbar{w}_j}^{\infty}
\Bigl[
w-\lbar{w}_j+\tilde V_m(w)-\tilde V_m(\lbar{w}_j)
\Bigr]\,dF(w)\\
&+
\frac{\beta(1-\gamma_{jj})}
{1-\beta\gamma_{mm}(1-\delta)}
\int_{\lbar{w}_m}^{\infty}
\Bigl[
w-\lbar{w}_m+\tilde V_s(w)-\tilde V_s(\lbar{w}_m)
\Bigr]\,dF(w)
\biggr]
\Biggr)
\Biggr\}
\end{align*}
As before, the search effort of the young worker is ambiguous with respect to that of the middle-aged worker, as is the effect on their reservation wage.
\medskip

%Turning to the evolution of the system more globally, we note that a worker is young for an average of $10$ years, middle-aged for $25$, and senior for $5$, hence participating in the labour force for an average of $40$ years. It follows that, at the stationary distribution, the proportion of workers in each age cohort must be $0.25$, $0.625$, and $0.125$ respectively.
Turning to the evolution of the system more globally, we note that per the representations of the age cohorts, a worker is in age cohort $g$ for duration $d_g$. It follows that, at the stationary distribution, we can define the proportion of workers in a given age cohort by:
\[
n_g\doteq\frac{d_g}{d_j+d_m+d_s}\quad\forall g\in\{j,m,s\}
\]
More generally, the expected length of time an agent stays in a certain age cohort is:
\[
\mathbb{E}[T_g]=\sum^\infty_{t=1}t\gamma_{gg}^{t-1}(1-\gamma_{gg})
\]
where we weight each time duration $t$ by the probability that they stay for precisely $t$, an and hence must leave in the next period. Solving the geometric series we obtain:
\[
\mathbb{E}[T_g]=\frac{1}{1-\gamma_{gg}}
\]
It follows that the transition probabilities are therefore given by:
\begin{comment}
\begin{align*}
10=\frac{1}{1-\gamma_{jj}},\quad 
25=\frac{1}{1-\gamma_{mm}},\quad 
5=\frac{1}{1-\gamma_{ss}}\\
\gamma_{jj}=0.9,\quad
\gamma_{mm}=0.96,\quad
\gamma_{ss}=0.8
\end{align*}
\end{comment}
\begin{align*}
d_j=\frac{1}{1-\gamma_{jj}},\quad 
d_m=\frac{1}{1-\gamma_{mm}},\quad 
d_s=\frac{1}{1-\gamma_{ss}}\\
\gamma_{jj}=\frac{d_j-1}{d_j},\quad
\gamma_{mm}=\frac{d_m-1}{d_m},\quad
\gamma_{ss}=\frac{d_s-1}{d_s}
\end{align*}
Putting everything together, let us denote the respective population weights of the unemployed within each cohort as $u_g$ for $g\in\{j,m,s\}$. Let $h_t$ denote our state vector of workers at $t$, hence given by:
\[
h_t \doteq
\bigl[(1-u_j),\, u_j,\, (1-u_m),\, u_m,\, (1-u_s),\, u_s\bigr]
\bigl(\operatorname{diag}(n_j,n_m,n_s)\otimes I_2\bigr)
\]
where each entry is weighted by the stationary population weight of its respective cohort. Furthermore, we note that a worker leaves unemployment, only if they both have an offer, and this offer is at a wage greater than their reservation wage in the period in which they can accept. Let us denote this by $p_g$ for each cohort such that:
\[
p_{gg'}\doteq\pi(a^*_g)(1-F(\lbar{w}_{g'}))
\]
Correspondingly, we denote the per-age labour-force transition matrix by:
\[
L_{gg'}
\doteq
\begin{bmatrix}
1-\delta & \delta \\
p_{gg'} & 1-p_{gg'}
\end{bmatrix},\quad
\mathcal L
\doteq
\begin{bmatrix}
L_{jj} & L_{jm} & \mathbf{0} \\
\mathbf{0} & L_{mm} & L_{ms} \\
\mathbf{0} & \mathbf{0} & L_{ss}
\end{bmatrix}.
\]
Hence, the full system evolves as:\footnote{The full matrix can be found in the appendix. In effect, the matrix is $6\times6$ with each odd row for employed workers at $t$ and each even row for unemployed workers at $t$, with the age cohort increasing every odd row, as one goes down the matrix.}
\begin{align*}
h_{t+1}&=h_t\bigr[\bigl((\Gamma\otimes\mathbf{1}_{2\times 2})\odot\mathcal{L}\bigr)+(1-\gamma_{ss})\bigl((e_s e_j')\otimes(\mathbf 1_2 e_u')\bigr)\bigl]\\
&=h_t\Pi
\end{align*}
where:\footnote{In particular, these elementary column vectors and the vector of $1$s combine to produce:
\[
\begin{bmatrix}
0 & 0 & 0 & 0 & 0 & 0 \\
0 & 0 & 0 & 0 & 0 & 0 \\
0 & 0 & 0 & 0 & 0 & 0 \\
0 & 0 & 0 & 0 & 0 & 0 \\
0 & 1 & 0 & 0 & 0 & 0 \\
0 & 1 & 0 & 0 & 0 & 0 \\
\end{bmatrix}
\]
which takes the last two entries of $h_t$ corresponding to the old age cohort (and hence the bottom two rows in the matrix), and sends them both to the unemployed young cohort at $t+1$, given by the second column.}
\[
e_j=
\begin{bmatrix}1\\0\\0\end{bmatrix},
\qquad
e_s=
\begin{bmatrix}0\\0\\1\end{bmatrix},
\qquad
e_u=
\begin{bmatrix}0\\1\end{bmatrix},
\qquad
\mathbf 1_2=
\begin{bmatrix}1\\1\end{bmatrix}.
\]
In effect, to obtain $h_{t+1}$, we multiply $H_t$ by a constructed transition matrix, which we denote as $\Pi$, where the first set of terms defines a $6\times6$ matrix comprised of the respective age and labour-state transitions, and the second set of terms sends the senior employed and unemployed workers into unemployed young workers with probability $1-\gamma_{ss}$. Since $\Pi$ is finite and row-stochastic, a stationary distribution $h^*$ exists such that:
\[
h^*=h^*\Pi
\]
In particular, under the usual positivity conditions ensuring irreducibility and aperiodicity, the system converges to it from any initial distribution. In effect, $h^*$ is the normalised row eigenvector of $\Pi$ with eigenvalue $1$.

Taking a more explicit analytic approach for this exposition, we get the following analytic formulae for the age-cohort-specific laws of motion, dictating the flows in and out of unemployment:
\begin{align*}
u_j'=&\;\gamma_{jj}\bigl((1-p_{jj})u_j+\delta(1-u_j))\bigr) + \frac{n_s}{n_j}(1-\gamma_{ss})\\
u_m'=&\;\gamma_{mm}\bigl((1-p_{mm})u_m+\delta(1-u_m))\bigr)\\
&+\frac{n_j}{n_m}(1-\gamma_{jj})\bigl((1-p_{jm})u_j+\delta(1-u_j))\bigr)\\
u'_s=&\;\gamma_{ss}\bigl((1-p_{ss})u_s+\delta(1-u_s))\bigr)\\
&+\frac{n_m}{n_s}(1-\gamma_{mm})\bigl((1-p_{ms})u_m+\delta(1-u_m))\bigr)\\
\end{align*}
Since these values are constant between periods in equilibrium, we can solve recursively starting at the unemployment rate of the young to get:
\begin{align*}
u_j^*&=\frac{\gamma_{jj}\delta+\frac{n_s}{n_j}(1-\gamma_{ss})}{1-\gamma_{jj}(1-\delta-p_{jj})}\\
u_m^*&=\frac{\gamma_{mm}\delta+\frac{n_j}{n_m}(1-\gamma_{jj})\left[\delta+(1-\delta-p_{jm})u_j^*\right]}{
1-\gamma_{mm}(1-\delta-p_{mm})}\\
u_s^*&=\frac{\gamma_{ss}\delta+\frac{n_m}{n_s}(1-\gamma_{mm})\left[
\delta+(1-\delta-p_{ms})u_m^*\right]}
{1-\gamma_{ss}(1-\delta-p_{ss})}
\end{align*}
Alternatively, solving in term of employment shares:
\begin{align*}
1-u_j'
=&\;
\gamma_{jj}
\left[
(1-\delta)(1-u_j)+p_{jj}u_j
\right],
\\
1-u_m'
=&\;
\gamma_{mm}
\left[
(1-\delta)(1-u_m)+p_{mm}u_m
\right]\\
&+\frac{n_j}{n_m}(1-\gamma_{jj})
\left[
(1-\delta)(1-u_j)+p_{jm}u_j
\right],
\\
1-u_s'
=&\;
\gamma_{ss}
\left[
(1-\delta)(1-u_s)+p_{ss}u_s
\right]\\
&+\frac{n_m}{n_s}(1-\gamma_{mm})
\left[
(1-\delta)(1-u_m)+p_{ms}u_m
\right].
\end{align*}
Solving at equilibrium, we have:
\begin{align*}
1-u_j^*
&=
\frac{\gamma_{jj}p_{jj}}
{1-\gamma_{jj}(1-\delta-p_{jj})}\\
1-u_m^*
&=
\frac{
\gamma_{mm}p_{mm}
+\frac{n_j}{n_m}(1-\gamma_{jj})\left[(1-\delta)+(p_{jm}-1+\delta)u_j^*\right]
}
{1-\gamma_{mm}(1-\delta-p_{mm})}\\
1-u_s^*
&=
\frac{
\gamma_{ss}p_{ss}
+\frac{n_m}{n_s}(1-\gamma_{mm})\left[(1-\delta)+(p_{ms}-1+\delta)u_m^*\right]
}
{1-\gamma_{ss}(1-\delta-p_{ss})}.
\end{align*}
In particular, we note that our two derivations provide the same values of $u^*_g$, $\forall g\in\{j,m,s\}$, since, per our age transition probabilities:\footnote{To see this even more clearly, note that $\frac{n_g}{n_g'}=\frac{d_g}{d_g'}$. Hence:
\[
    \frac{n_s}{n_j}=\frac{1-\gamma_{jj}}{1-\gamma_{ss}},\quad
    \frac{n_j}{n_m}=\frac{1-\gamma_{mm}}{1-\gamma_{jj}},\quad
    \frac{n_m}{n_s}=\frac{1-\gamma_{ss}}{1-\gamma_{mm}}
\]
}
\[
\frac{n_s}{n_j}(1-\gamma_{ss})
=
1-\gamma_{jj},
\qquad
\frac{n_j}{n_m}(1-\gamma_{jj})
=
1-\gamma_{mm},
\qquad
\frac{n_m}{n_s}(1-\gamma_{mm})
=
1-\gamma_{ss}.
\]
\medskip\\
\noindent Before proceeding to the numerical methods section, we also deduce an expression for $G(w)$, the distribution of wages paid. In particular, since for a given age cohort $g\in\{j,m,s\}$, wages are only realised upon acceptance, the wage distribution coincides with the distribution of accepted offers, which is in turn a left-truncated distribution of $F(w)$. We note that:
\[
\mathbb{P}(W\leq w| W \geq \lbar{w}_g)=\frac{\mathbb{P}(\lbar{w}_g\leq W\leq w)}{\mathbb{P}(W\geq \lbar{w}_g)}=\frac{F(w)-F(\lbar{w}_g)}{1-F(\lbar{w}_g)}
\]
So, the corresponding wage distribution for $g$ is:
\[
G_g(w)=
\begin{cases}
    0,\quad &w<\lbar{w}_g\\
    \frac{F(w)-F(\lbar{w}_g)}{1-F(\lbar{w}_g)},
    \quad &w\geq\lbar{w}_g
\end{cases}
\]
Finally, aggregating across cohorts, the distribution of wages among employed workers is:
\[
G(w)=\frac{\sum_{g\in\{j,m,s\}} n_gG_g(w)}{\sum_{g\in\{j,m,s\} n_g(1-u_g^*)}}
\]
\section{Numerical Solution}
Given our expressions we solve the equilibrium recursively, starting at the senior problem. In our particular case, we first define a set of bounds for $\phi$, starting with an initial intermediate guess, and use a root-finding algorithm to produce the next guess of $\phi$, after having recursively having solved for the reservation wages.

To begin, we first infer that the each $t$ corresponds to a month. Hence, to compute the values of $\gamma_{gg}$, $\forall g\in\{j,m,s\}$. Since an individual is on average young for 10 years, middle-aged for 25, and senior for 5, we can multiply each of these values by 12 and solve to get:
\[
\gamma_{jj}=\frac{119}{120},\quad\gamma_{mm}=\frac{299}{300},\quad\gamma_{ss}=\frac{59}{60}
\]
We note that once a middle-aged worker is unemployed, the duration spent in the middle-aged unemployed state ends either because they find an acceptable job while remaining middle-aged, or because they age into the senior cohort. Hence, the per-period exit probability from middle-aged unemployment is:
\[
\gamma_{mm}\pi(a^*_m)(1-F(\lbar{w}_m))+(1-\gamma_{mm})
\]
using this interpretation for our calibration, it follows from the geometric distribution that the average duration spent unemployed while middle-aged is:
\[
13=\frac{1}{\gamma_{mm}\pi(a^*_m)(1-F(\lbar{w}_m))+(1-\gamma_{mm})}
\]
Since $a_m^*$ (and $\lbar{w}_m$) depend on $\phi$, we solve for $\phi$ by treating the above condition as a scalar equation. As an initialisation, we take a small positive lower bound for $\phi$ and construct an upper bound via geometric expansion, starting at $1$. For the purposes of this solution, we restrict our search to a unique $\phi$,\footnote{The uniqueness of $\phi$ can only be guaranteed by numerical solution. In the appendix, we plot our root-finding residual on $\phi$ over a large range of values, strongly suggesting a unique solution.} which we find by first expanding our upper-bracket geometrically until the right-hand-side of the above equation provides a value below 13. After this, we solve on our interval using Brent's method.\footnote{This relies upon the implicit assumption that the number of iterations to find the crossing point is small, such that the time saved using the super-linear convergence of Brent outweighs Picard iteration. For a discussion on why Brent is chosen over safeguarded Newton-Raphson, see the appendix.}

For each candidate value of $\phi$, we solve the reservation wages using a safeguarded Newton-Raphson method.\footnote{For a longer discussion on our choice of solution method, and a broader explanation of how it is implemented, see the appendix.} Since the average middle-aged unemployment duration depends only on $\lbar{w}_s$ and $\lbar{w}_m$, we solve for these sequentially:\footnote{To make the code solve in a more streamlined manner, some of the fixed-point formulae were re-arranged, please check the appendix to see these versions.} first $\lbar{w}_s$, and then $\lbar{w}_m$, reflecting the dependence of the latter on the former. 

Given these values, we compute the implied average middle-aged unemployment duration and evaluate the calibration residual. The value of $\phi$ is then updated using Brent's method.

Once we obtain a value of $\phi$ that brings the calibration residual within a prescribed tolerance, we solve for the remaining equilibrium objects, in particular the full set of reservation wages and the steady-state employment shares, which are then used to construct $G(w)$. 

\subsection{Age, Research Effort, and Reservation Wage}
Solving at the given parameter values, we obtain a calibrated $\phi$ value of around $0.00055$. At equilibrium, we have the following reservation wages and search efforts for each age cohort:
\begin{table}[!ht]
\centering
\begin{tabular}{lcc}
\hline
Age cohort & Reservation wage & Search effort \\
\hline
Young & 56.9672 & 162.155 \\
Middle-aged & 55.4289 & 143.874 \\
Senior & 50 & 0 \\
\hline
\end{tabular}
\caption{Calibrated reservation wages and search efforts by age cohort.}
\label{tab:rwage_search_effort}
\end{table}

The search-effort ordering in Table~\ref{tab:rwage_search_effort} can be understood from the closed-form effort rules. Let $\mathcal G_g$ denote the expected gain term entering the effort choice of cohort $g$, so that:
\[
a_g^*
=
\max\left\{
0,\frac{1}{\phi}
\log\left(\frac{\phi\mathcal G_g}{\psi}\right)
\right\}.
\]
Since the logarithm is increasing, search effort is increasing in $\mathcal G_g$ whenever the interior condition holds. The affine continuation-value terms help determine these gain terms. In particular, the middle-aged continuation integral can be written as:\footnote{See appendix for the derivation.}
\[
\int_{\underline w_m}^{\infty}
\left[
w-\underline w_m+\tilde V_s(w)-\tilde V_s(\underline w_m)
\right]\,dF(w)
=
(1+B_s)
\int_{\underline w_m}^{\infty}
(w-\underline w_m)\,dF(w),
\]
where:
\[
B_s
=
\frac{\beta(1-\gamma_{mm})(1-\delta)}
{1-\beta\gamma_{ss}(1-\delta)}
>0.
\]
Similarly, the young surplus integral is scaled by $1+B_m$, where:
\[
B_m
=
\frac{\beta(1-\gamma_{jj})(1-\delta)(1+B_s)}
{1-\beta\gamma_{mm}(1-\delta)}
>0.
\]
Thus, a marginally higher wage offer is more valuable for younger workers because it raises not only current employment value, but also continuation values in later age states. This creates a force toward larger gain terms $\mathcal G_g$, and therefore higher search effort, for younger cohorts. In the calibrated solution, this force is strong enough that young workers choose the highest search effort, followed by middle-aged workers, while seniors remain at the corner.

The reservation-wage ordering should be interpreted more cautiously. Higher search effort does not mechanically imply a higher reservation wage, because the reservation-wage fixed point also contains the flow search cost $-\psi a_g^*$ and the reservation wage itself changes the tail integrals. Hence the ranking of reservation wages is determined by the full fixed-point system rather than by the effort ordering alone. Numerically, however, the calibrated solution gives the same age ranking for reservation wages as for search efforts.

For seniors, the calibrated solution is at a corner. Using our closed form expression for $a^*_s$ from our analytic section, the senior search rule implies positive search only if:
\[
\frac{\phi\tilde\beta_sR(\underline w_s)}{\psi}>1,
\qquad
\tilde\beta_s
\doteq
\frac{\beta\gamma_{ss}}
{1-\beta\gamma_{ss}(1-\delta)}.
\]
At the calibrated value of $\phi$ implied by the 13-month middle-aged unemployment-duration target, this condition fails at $\underline w_s=b$. This is despite the fact that $\tilde\beta_s$ is itself fairly large, reflecting the high values of $\beta$ and $\gamma_{ss}$ together with the relatively low job-destruction rate $\delta$. The corner therefore does not arise because senior employment spells are mechanically worthless; rather, at the calibrated value of $\phi$, the product of search efficiency and expected acceptable-offer surplus is still too small relative to the flow search cost $\psi$. Therefore $a_s^*=0$, and since $\pi(0)=0$, the senior fixed-point equation collapses to $\underline w_s=b$. This corner is not a general feature of the model: if the target duration is shortened enough, the calibrated value of $\phi$ rises, and senior search can become interior. Numerically, this occurs around a 7-month target in the present parametrisation. Thus the zero senior search effort in Table~\ref{tab:rwage_search_effort} should be interpreted as a consequence of the particular calibration target rather than as a mechanical property of the senior problem.

\subsection{Unemployment Rates and $G(w)$}
Table~\ref{tab:unemployment_by_age} reports two related unemployment measures: the unemployment rate within each age cohort, and each cohort's unemployed workers as a share of the total population.
\begin{table}[!ht]
\centering
\begin{tabular}{lccc}
\hline
$g$ & $n_g$ & $u^*_g$ & $n_g u^*_g$ \\
\hline
$j$ & 0.25 & 0.233388 & 0.058347 \\
$m$ & 0.625 & 0.185678 & 0.116049 \\
$s$ & 0.125 & 0.589244 & 0.073655 \\
\hline
Total & 1 & 0.248051 & 0.248051 \\
\hline
\end{tabular}
\caption{Steady-state unemployment proportions by age cohort.}
\label{tab:unemployment_by_age}
\end{table}
The senior cohort has the highest within-cohort unemployment rate, reflecting the corner solution for senior search effort. However, because seniors make up only $12.5\%$ of the population, their contribution to aggregate unemployment is smaller than that of middle-aged workers.

We first note that the model does not include a firm side, and hence there is no scope for congestion externalities from workers competing for vacancies. The only labour-market choice variable is individual search effort. Unemployment differences across cohorts therefore arise from the interaction of search effort, reservation wages, and age-transition probabilities.

This can be seen from the unemployment-to-employment transition probability:
\[
p_{gg'}=\pi(a_g^*)\left[1-F(\underline w_{g'})\right].
\]
A cohort has a lower unemployment rate when its workers search more intensely, or when the relevant next-period reservation wage is lower, raising the probability that a wage draw is accepted. Age transitions also matter because unemployed workers may enter the next age cohort before finding a job.

The unemployment rates are not monotone in age: middle-aged workers have a lower unemployment rate than young workers, while seniors have the highest unemployment rate. The high senior unemployment rate reflects their low search effort and, in the calibrated solution, the corner \(a_s^*=0\). The relatively high young unemployment rate has a different source: all new labour-force entrants enter as unemployed young workers. Thus the young cohort receives a direct inflow into unemployment each period from the replacement of retiring seniors. Middle-aged workers do not receive an analogous entry flow, which helps explain why their unemployment rate is lower despite having slightly lower search effort than young workers.

These employment shares also determine the weights in the aggregate accepted-wage distribution $G(w)$ derived in the previous section: each cohort's truncated wage distribution is weighted by its employment mass. Below, we plot both $G(w)$ and $F(w)$ starting at $50$, the truncation point under this parametrisation, representing the lowest reservation wage rate, or, in this specific case, given senior workers are at a corner solution, the benefit rate $b$.
\begin{figure}[!ht]
\centering
\includegraphics[width=0.8\textwidth]{\detokenize{MATLAB code/wage_distribution_plot.png}}
\caption{Offer distribution $F(w)$ and aggregate accepted-wage distribution $G(w)$, plotted from the lowest calibrated reservation wage.}
\label{fig:wage_distribution}
\end{figure}
As we can see, the accepted-wage distribution $G(w)$ first-order stochastically dominates the offer distribution $F(w)$. This reflects selection into employment: wages are only observed when offers exceed the relevant reservation wage, so accepted wages are drawn from the upper tails of the offer distribution.

As for the shape of the distribution itself, although each cohort-specific accepted-wage distribution is simply a left-truncated version of the offer distribution, the aggregate distribution $G(w)$ is a mixture of these cohort-specific distributions. Since cohorts have different reservation wages, the components of the mixture enter at different wage thresholds. As $w$ crosses each reservation wage, an additional cohort begins contributing to the aggregate distribution. This can generate visible kinks, shoulders, or bumps in $G(w)$, especially when the reservation wages are sufficiently separated or when a cohort has a large employment weight.

\appendix
\section{The Transition Matrix}

Here we write out more fully the transition matrix used to compute the stationary distribution. The state vector is:
\[
h_t
\doteq
\bigl[
n_j(1-u_j),\, n_j u_j,\,
n_m(1-u_m),\, n_m u_m,\,
n_s(1-u_s),\, n_s u_s
\bigr],
\]
where the first entry in each age pair is employment and the second is unemployment. Recall that:
\[
p_{gg'}=\pi(a_g^*)\left[1-F(\underline w_{g'})\right],
\]
so \(p_{gg'}\) is the probability that an unemployed worker currently of age \(g\) receives and accepts an offer when next period's age is \(g'\). For each feasible age transition \(g\to g'\), the labour-market transition block is:
\[
L_{gg'}
=
\begin{bmatrix}
1-\delta & \delta\\
p_{gg'} & 1-p_{gg'}
\end{bmatrix}.
\]
For example, if a worker is currently unemployed and moves from \(j\) to \(m\), the relevant job-finding probability is \(p_{jm}\), because the acceptance threshold is the middle-aged reservation wage \(\underline w_m\).

The age transition component is:
\[
\Gamma\otimes \mathbf 1_{2\times2}
=
\begin{bmatrix}
\gamma_{jj} & \gamma_{jj} & 1-\gamma_{jj} & 1-\gamma_{jj} & 0 & 0\\
\gamma_{jj} & \gamma_{jj} & 1-\gamma_{jj} & 1-\gamma_{jj} & 0 & 0\\
0 & 0 & \gamma_{mm} & \gamma_{mm} & 1-\gamma_{mm} & 1-\gamma_{mm}\\
0 & 0 & \gamma_{mm} & \gamma_{mm} & 1-\gamma_{mm} & 1-\gamma_{mm}\\
0 & 0 & 0 & 0 & \gamma_{ss} & \gamma_{ss}\\
0 & 0 & 0 & 0 & \gamma_{ss} & \gamma_{ss}
\end{bmatrix},
\]
while the corresponding labour-market block matrix is:
\[
\mathcal L
=
\begin{bmatrix}
L_{jj} & L_{jm} & \mathbf 0\\
\mathbf 0 & L_{mm} & L_{ms}\\
\mathbf 0 & \mathbf 0 & L_{ss}
\end{bmatrix}.
\]
The Hadamard product \((\Gamma\otimes \mathbf 1_{2\times2})\odot\mathcal L\) applies the appropriate age-transition probability to each feasible labour-market transition. Finally, retiring senior workers are replaced by unemployed young workers. This is captured by:
\[
(1-\gamma_{ss})
\begin{bmatrix}
0 & 0 & 0 & 0 & 0 & 0\\
0 & 0 & 0 & 0 & 0 & 0\\
0 & 0 & 0 & 0 & 0 & 0\\
0 & 0 & 0 & 0 & 0 & 0\\
0 & 1 & 0 & 0 & 0 & 0\\
0 & 1 & 0 & 0 & 0 & 0
\end{bmatrix}.
\]
Putting these pieces together, the full row-stochastic transition matrix is:
\[
\resizebox{\textwidth}{!}{$
\Pi
=
\begin{bmatrix}
\gamma_{jj}(1-\delta) & \gamma_{jj}\delta
& (1-\gamma_{jj})(1-\delta) & (1-\gamma_{jj})\delta
& 0 & 0\\
\gamma_{jj}p_{jj} & \gamma_{jj}(1-p_{jj})
& (1-\gamma_{jj})p_{jm} & (1-\gamma_{jj})(1-p_{jm})
& 0 & 0\\
0 & 0
& \gamma_{mm}(1-\delta) & \gamma_{mm}\delta
& (1-\gamma_{mm})(1-\delta) & (1-\gamma_{mm})\delta\\
0 & 0
& \gamma_{mm}p_{mm} & \gamma_{mm}(1-p_{mm})
& (1-\gamma_{mm})p_{ms} & (1-\gamma_{mm})(1-p_{ms})\\
0 & 1-\gamma_{ss}
& 0 & 0
& \gamma_{ss}(1-\delta) & \gamma_{ss}\delta\\
0 & 1-\gamma_{ss}
& 0 & 0
& \gamma_{ss}p_{ss} & \gamma_{ss}(1-p_{ss})
\end{bmatrix}
$}.
\]
The stationary distribution \(h^*\) is therefore the normalized row vector satisfying:
\[
h^*=h^*\Pi,
\qquad
\sum_i h_i^*=1.
\]

\section{The Choice of Solution Method}
\subsection{Solving for $\lbar{w}_g$, $\forall g\in\{j,m,s\}$}
We begin by noting that none of the wage fixed-point equations are contractions. In particular, let:
\begin{align*}
T_s(x)
=
\max_{a_s\geq0}
\left\{
b-\psi a_s+
\tilde\beta_s\pi(a_s)R(x)
\right\},\\
R(x)\doteq\sigma\varphi\biggl(\frac{x-\mu}{\sigma}\biggr) + (\mu - x)\Biggl(1 - \Phi\biggl(\frac{x-\mu}{\sigma}\biggr)\Biggr),
\qquad
\tilde\beta_s\doteq
\frac{\beta\gamma_{ss}}
{1-\beta\gamma_{ss}(1-\delta)}.
\end{align*}
Where, for any $x,y\in\supp(F)$,
\[
|T_s(x)-T_s(y)|
\leq
\tilde\beta_s |R(x)-R(y)|.
\]
Since $R'(x)=-(1-F(x))$ and $0\leq 1-F(x)\leq1$, the mean value theorem gives
\[
|R(x)-R(y)|\leq |x-y|.
\]
Hence,
\[
|T_s(x)-T_s(y)|
\leq
\tilde\beta_s |x-y|.
\]
Therefore, since $\tilde\beta_s>1$, $T_s$ is not a contraction (on $\supp(F)$). A similar logic follows for the remaining two fixed-point equations.

Nonetheless, if one still wishes to use solution methods which maintain contractive properties, it remains possible to directly solve from the Bellman equations via VFI. Alternatively, we could pursue an endogenous cut-off formulation to eliminate the computationally-intensive acceptance maximisation and solve the model via value function iteration. However, this approach still requires repeated updates of the value functions and associated integrals, leading to at best linear convergence. By contrast, directly solving the reservation-wage equation via a safeguarded Newton method reduces the problem to a scalar root-finding exercise with quadratic local convergence, and is therefore computationally more efficient.

In practical terms, this amounts to searching for the unique root of:
\[
\Psi_g(\lbar{w}_g)=\lbar{w}_g-T_g(\lbar{w}_g).
\]
Taking a starting interval, for instance \((b,\mu+2\sigma)\), and verifying that the sign does indeed change over this interval, we maintain a bracketing interval \([a_k,b_k]\) such that
\[
\Psi_g(a_k)\Psi_g(b_k)<0
\]
at each iteration.

At a given iterate \(x_k\in[a_k,b_k]\), the method first computes the Newton candidate:
\[
x_{k+1}^{N}
=
x_k-\frac{\Psi_g(x_k)}{\Psi_g'(x_k)}.
\]
If this Newton candidate remains inside the current bracket and delivers sufficient progress (e.g. reduces the residual or lies sufficiently close to the current iterate), we accept it:
\[
x_{k+1}=x_{k+1}^{N}.
\]
Otherwise, we reject the Newton step and instead take the bisection update:
\[
x_{k+1}
=
\frac{a_k+b_k}{2}.
\]
After evaluating \(\Psi_g(x_{k+1})\), the bracket is updated according to the sign of the residual:
\[
[a_{k+1},b_{k+1}]
=
\begin{cases}
[x_{k+1},b_k],
& \text{if } \Psi_g(x_{k+1})\Psi_g(a_k)>0,\\
[a_k,x_{k+1}],
& \text{if } \Psi_g(x_{k+1})\Psi_g(a_k)<0.
\end{cases}
\]
This guarantees that the root remains bracketed throughout the iteration.

The reason this method is particularly well suited to the present problem is that \(\Psi_g'(\cdot)\) can be computed analytically. Hence, when the current iterate is sufficiently close to the root and \(\Psi_g'(\lbar{w}_g)\neq 0\), the Newton step inherits the usual local convergence properties of Newton-Raphson. To see this, a Taylor expansion around the true root \(\lbar{w}_g^*\) gives:
\[
0
=
\Psi_g(\lbar{w}_g^*)
=
\Psi_g(x_k)
+
\Psi_g'(x_k)(\lbar{w}_g^*-x_k)
+
O\bigl((\lbar{w}_g^*-x_k)^2\bigr),
\]
so that the Newton update eliminates the first-order error term. Consequently, once the iterates enter a neighbourhood of the root, convergence is quadratic under the usual regularity conditions. Away from the root, the bisection safeguard preserves global robustness by preventing the algorithm from leaving the bracketing interval.

This hybrid strategy therefore combines the reliability of bisection with the speed of Newton-Raphson. Since \(\Psi_g(\cdot)\) is continuous and strictly increasing, the maintained bracket contains a unique root, while the safeguarded Newton updates provide an efficient and reliable method for computing the reservation wage for each cohort. This is also why we do not use Picard iteration for the reservation-wage equations. Iterating \(\underline w_g^{k+1}=T_g(\underline w_g^k)\) would deliver at best linear convergence when it converges, while the wage equations have analytical derivatives that allow the safeguarded Newton step to obtain quadratic local convergence.

\subsection{Solving for $\phi$}
For the outer calibration, define the residual:
\[
\Omega(\phi)
\doteq
\frac{1}{
\gamma_{mm}\pi(a_m^*(\phi),\phi)
\left[1-F(\underline w_m(\phi))\right]
+(1-\gamma_{mm})
}
-13.
\]
The calibrated value of \(\phi\) solves \(\Omega(\phi)=0\).

We solve for \(\phi\) using \texttt{fzero}, native to \texttt{MATLAB}. This routine implements a Brent-type method, combining bisection with interpolation steps such as the secant method and inverse quadratic interpolation. We choose a bracketing interval \([\phi_L,\phi_H]\) such that \(\Omega(\phi_L)\Omega(\phi_H)<0\). The bisection component preserves the bracketing of the root, while the interpolation steps typically accelerate convergence relative to pure bisection.

In the secant case, given two previous residual evaluations at distinct values of \(\phi\), the proposed update is:
\[
\phi_{k+1}^{S}
=
\phi_k
-
\Omega(\phi_k)
\frac{\phi_k-\phi_{k-1}}{\Omega(\phi_k)-\Omega(\phi_{k-1})}.
\]
When three previous residual evaluations are available at distinct values of \(\phi\), and their residual values are also distinct, the method may instead employ inverse quadratic interpolation. Let:
\[
y_i\doteq \Omega(\phi_i),
\qquad i\in\{0,1,2\}.
\]
The method fits a quadratic approximation to the inverse map \(\phi=\phi(y)\) through the three points \((y_i,\phi_i)\), and evaluates this approximation at \(y=0\). Using the Lagrange form of the interpolating polynomial, the proposed inverse-quadratic update is:
\[
\phi_{k+1}^{Q}
=
\phi_0
\frac{(0-y_1)(0-y_2)}{(y_0-y_1)(y_0-y_2)}
+
\phi_1
\frac{(0-y_0)(0-y_2)}{(y_1-y_0)(y_1-y_2)}
+
\phi_2
\frac{(0-y_0)(0-y_1)}{(y_2-y_0)(y_2-y_1)}.
\]
Equivalently, this is the value of \(\phi\) implied by the fitted inverse quadratic when the residual is zero. This yields a higher-order approximation to the root without requiring a derivative of \(\Omega(\phi)\).

The interpolation step is not accepted automatically. Brent-type methods use an interpolation step only when the required distinct function values are available and the proposed point satisfies safeguard conditions: it must remain inside the current bracket and deliver sufficient progress relative to a bisection step. If these conditions fail, the method falls back to bisection:
\[
\phi_{k+1}=\frac{\phi_L+\phi_H}{2},
\]
thereby preserving the bracketing condition at all times.

A natural alternative would be to rewrite the calibration equation as a fixed-point problem, \(\phi=G(\phi)\), and use Picard iteration:
\[
\phi_{k+1}=G(\phi_k).
\]
However, even when this iteration converges, its convergence is generally only linear. To see this, suppose \(\phi^*=G(\phi^*)\) and \(G\) is differentiable in a neighbourhood of \(\phi^*\). A first-order Taylor expansion gives:
\[
\phi_{k+1}-\phi^*
=
G'(\phi^*)(\phi_k-\phi^*)+O\left((\phi_k-\phi^*)^2\right).
\]
Thus, unless \(G'(\phi^*)=0\), the error shrinks only proportionally to the previous error. By contrast, a Brent-type method keeps the robustness of bisection but can exploit secant or inverse-quadratic interpolation steps when these are reliable. Under the usual smoothness conditions, secant-type updates converge superlinearly, while inverse quadratic interpolation has an asymptotic order above linear as well. These interpolation steps use local information about \(\Omega(\phi)\) to jump closer to the root, which is why \texttt{fzero} is typically faster than Picard iteration for a scalar calibration equation, while still falling back to bisection when the interpolation step is unsafe.

This hybrid strategy combines the global robustness of bisection with the faster local convergence of interpolation-based methods. While safeguarded Newton-Raphson is effective for the reservation-wage equations due to the availability of analytical derivatives, applying Newton-Raphson to the outer calibration problem is less attractive. The residual \(\Omega(\phi)\) depends on \(a_m^*(\phi)\), \(\underline w_m(\phi)\), and \(\underline w_s(\phi)\), where the reservation wages are themselves obtained by solving nonlinear equations for each candidate value of \(\phi\). Moreover, optimal search effort is defined by a corner condition,
\[
a_g^*=\max\{0,\cdots\},
\]
so the derivative of the outer residual is piecewise and may be nonsmooth at values of \(\phi\) where the search-effort constraint binds or ceases to bind. Differentiating \(\Omega(\phi)\) is therefore possible in principle, but would require differentiating through the nested equilibrium solver while accounting for these corner regimes. MATLAB's \texttt{fzero} avoids this additional complexity because it only requires evaluations of \(\Omega(\phi)\), while its bisection safeguard preserves the root bracket even when interpolation steps are unreliable.

By contrast, for the reservation-wage equations we use safeguarded Newton-Raphson because the relevant derivatives are available analytically from the envelope theorem. Inverse quadratic interpolation is most useful when these derivatives are unavailable or inconvenient to compute; using it for the inner problem would therefore add little relative to Newton's quadratic local convergence.

\subsection{Numerical Check of the Calibration Residual}

Since the calibration residual depends on recursively solved reservation wages and piecewise-defined search efforts, we do not provide an analytic proof of global uniqueness for the root in \(\phi\). Instead, we numerically evaluate \(\Omega(\phi)\) over a wide grid of positive values of \(\phi\). Figure~\ref{fig:calibration_residual} plots the resulting residual. For small values of \(\phi\), search is effectively inactive, so the exit probability from middle-aged unemployment is close to the aging probability \(1-\gamma_{mm}\). The implied duration is therefore close to:
\[
\frac{1}{1-\gamma_{mm}}=300
\]
months, giving the high plateau on the left of the plot. As \(\phi\) increases, middle-aged search becomes productive and the residual falls sharply through zero. For larger values of \(\phi\), the residual remains negative and flattens, reflecting the endogenous adjustment of search effort and reservation wages as search efficiency rises. Hence, over the grid considered, the residual is positive on the low-\(\phi\) plateau, crosses zero once during the sharp transition, and remains negative thereafter. This provides numerical evidence for a unique calibrated root on the relevant search interval and motivates the bracketing interval used by \texttt{fzero}.

\begin{figure}[!ht]
\centering
\includegraphics[width=0.8\textwidth]{\detokenize{MATLAB code/calibration_residual_plot.png}}
\caption{Calibration residual \(\Omega(\phi)\) over a log-spaced grid of values for \(\phi\).}
\label{fig:calibration_residual}
\end{figure}

\section{Simplifying $T_g(\lbar{w}_g\in\{j,m\})$ for the Solver}
While $R(x)$ remains an integral that lends itself well to the \texttt{MATLAB} normal pdf and cdf functions, the remaining integrals in the middle-aged and young reservation wage equations can do with more simplification, so as to prevent the evaluation of ad hoc integrals and save on computational power.

\subsection{Evaluation of the Middle-Aged Surplus Integral}

For the numerical implementation, it is useful to simplify the integral appearing in the middle-aged reservation-wage equation:
\[
I_m(x)
\doteq
\int_x^\infty
\Bigl[
w-x+\tilde V_s(w)-\tilde V_s(x)
\Bigr]\,dF(w).
\]
Where we recall that:
\[
\tilde V_s(w)=
\beta(1-\gamma_{mm})
\left[
(1-\delta)V_{e,s}(w)+\delta V_{u,s}
\right],\quad
V_{e,s}(w)=
\frac{
w+\beta\gamma_{ss}\delta V_{u,s}
}{
1-\beta\gamma_{ss}(1-\delta)
}.
\]
In effect, $V_{e,s}(w)$ is affine in $w$, and therefore so is
$\tilde{V}_s(w)$.
\smallskip\\
\noindent From here, let us define $\tilde{V}_s(w)$ in terms of $A_s$ and $B_s$:
\[
\tilde V_s(w)=A_s+B_s w,
\]
where:
\[
B_s
\doteq
\frac{\beta(1-\gamma_{mm})(1-\delta)}{1-\beta\gamma_{ss}(1-\delta)},\quad
A_s
\doteq
\beta(1-\gamma_{mm})
\left[
\frac{(1-\delta)\beta\gamma_{ss}\delta V_{u,s}}{1-\beta\gamma_{ss}(1-\delta)}
+\delta V_{u,s}
\right].
\]
Further recall that from the senior reservation-wage condition:
\[
V_{u,s}
=
\frac{\underline w_s}{1-\beta\gamma_{ss}},
\]
hence both \(A_s\) and \(B_s\) are known once \(\underline w_s\) has been solved.
\smallskip\\
\noindent Substituting \(\tilde V_s(w)=A_s+B_s w\) into \(I_m(x)\), we obtain:
\begin{align*}
I_m(x)
&=
\int_x^\infty
\Bigl[
w-x+A_s+B_s w-(A_s+B_s x)
\Bigr]\,dF(w)\\
&=
\int_x^\infty
\Bigl[
(1+B_s)(w-x)
\Bigr]\,dF(w)\\
&=
(1+B_s)
\int_x^\infty (w-x)\,dF(w).
\end{align*}
It follows that:
\[
I_m(x)
=
(1+B_s)R(x), \quad R(x)
\doteq
\int_x^\infty (w-x)\,dF(w)
\]
Thus, the middle-aged surplus integral does not require numerical quadrature. Since $\tilde V_s(\cdot)$ is affine in the wage, the integral can be evaluated directly from the same truncated-surplus function $R(x)$ used in the senior problem.

\subsection{Evaluation of the Young Surplus Integral}

The young worker's reservation-wage equation can be simplified in the same way. Define:
\[
I_j(x)
\doteq
\int_x^\infty
\Bigl[
w-x+\tilde V_m(w)-\tilde V_m(x)
\Bigr]\,dF(w).
\]
To evaluate this term, first note that the middle-aged employment value is affine in the wage once $\underline w_s$ and $\underline w_m$ have been solved. From the reservation-wage condition,
\[
\underline w_m=(1-\beta\gamma_{mm})V_{u,m}-\tilde V_s(\underline w_m),
\]
we can recover:
\[
V_{u,m}
=
\frac{\underline w_m+\tilde V_s(\underline w_m)}
{1-\beta\gamma_{mm}}.
\]
Since:
\[
V_{e,m}(w)
=
\frac{
w+\beta\gamma_{mm}\delta V_{u,m}+\tilde V_s(w)
}{
1-\beta\gamma_{mm}(1-\delta)
},
\]
and $\tilde V_s(w)=A_s+B_s w$, it follows that $V_{e,m}(w)$ is affine in $w$. Consequently:
\[
\tilde V_m(w)
\doteq
\beta(1-\gamma_{jj})
\left[
(1-\delta)V_{e,m}(w)+\delta V_{u,m}
\right]
=A_m+B_m w,
\]
where the slope is:
\[
B_m
\doteq
\frac{\beta(1-\gamma_{jj})(1-\delta)(1+B_s)}
{1-\beta\gamma_{mm}(1-\delta)}.
\]
Substituting $\tilde V_m(w)=A_m+B_m w$ into $I_j(x)$ gives:
\begin{align*}
I_j(x)
&=
\int_x^\infty
\Bigl[
w-x+A_m+B_mw-(A_m+B_mx)
\Bigr]\,dF(w)\\
&=
(1+B_m)\int_x^\infty(w-x)\,dF(w)\\
&=(1+B_m)R(x).
\end{align*}
Hence the young worker's own-age surplus integral can also be evaluated without numerical quadrature.

The second integral in the young worker's fixed-point equation corresponds to the continuation surplus from entering the middle-aged cohort. Using the simplification above for the middle-aged surplus integral, this term is:
\[
\int_{\underline w_m}^{\infty}
\Bigl[
w-\underline w_m+\tilde V_s(w)-\tilde V_s(\underline w_m)
\Bigr]\,dF(w)
=
(1+B_s)R(\underline w_m).
\]
Thus, after solving $\underline w_s$ and $\underline w_m$, the young worker's gain term can be written entirely in terms of $R(\cdot)$ and known continuation-value slopes. This is the expression used in \texttt{solve\_wj.m}.

\end{document}
